\clearpage
\section{연습문제}
\label{sec:cyclotomic-cyclic-exercises}

문제는 A, B, C의 세 단계로 나뉜다. 별표 문제는 다음 두 절에 상세 풀이가 있다.

\subsection*{A. 계산과 구체적 예제}

\begin{exercise}[*]
$\mu_{12}(\C)$의 모든 원소를 $\zeta_{12}$의 거듭제곱으로 적고, 그중 원시 $12$차 단위근을 모두 골라라. 각 원소의 위수도 구하라.
\end{exercise}

\begin{exercise}
다음을 계산하라.
\[
  \varphi(72),\qquad
  \varphi(100),\qquad
  \varphi(105),\qquad
  \sum_{d\mid20}\varphi(d).
\]
\end{exercise}

\begin{exercise}[*]
약수별 분해와 소수붙이기 공식을 사용하여 다음 원분다항식을 계산하라.
\[
  \Phi_8,\quad\Phi_9,\quad\Phi_{10},
  \quad\Phi_{12},\quad\Phi_{15},\quad\Phi_{18}.
\]
\end{exercise}

\begin{exercise}
$X^{12}-1$을 $\Q[X]$에서 기약인수들의 곱으로 완전히 분해하라. 각 인수의 근이 갖는 정확한 위수도 적어라.
\end{exercise}

\begin{exercise}[*]
$L=\Q(\zeta_8)$에 대해 다음을 수행하라.
\begin{enumerate}[label=(\alph*)]
  \item $\Gal(L/\Q)$를 $(\Z/8\Z)^\times$와 동일시하라.
  \item $L=\Q(i,\sqrt2)$임을 보여라.
  \item 세 위수 $2$ 부분군과 세 이차 중간체
  \[
    \Q(i),\quad\Q(\sqrt2),\quad\Q(\sqrt{-2})
  \]
  를 정확히 대응시켜라.
\end{enumerate}
\end{exercise}

\begin{exercise}[*]
$t=\zeta_5+\zeta_5^{-1}$라 하자. $t$의 최소다항식을 구하고
\[
  \Q(\zeta_5)^+=\Q(\sqrt5)
\]
임을 보여라. 이어서 $\Q(\zeta_5)/\Q$의 중간체 격자를 그려라.
\end{exercise}

\begin{exercise}[*]
$t=\zeta_7+\zeta_7^{-1}$라 하자. 재귀
\[
  u_{k+1}=tu_k-u_{k-1},
  \qquad u_k=\zeta_7^k+\zeta_7^{-k}
\]
를 사용하여 $t$의 최소다항식이
\[
  X^3+X^2-2X-1
\]
임을 증명하라.
\end{exercise}

\begin{exercise}
$\Q(\zeta_{12})=\Q(i,\sqrt3)$임을 보이고 갈루아군과 세 이차 중간체를 모두 구하라.
\end{exercise}

\begin{exercise}[*]
$\Phi_7$을 $\F_2[X]$에서 인수분해하고 각 인수의 기약성을 확인하라. 원시 $7$차 단위근의 분해체와 갈루아군도 구하라.
\end{exercise}

\begin{exercise}
다음을 계산하라.
\begin{enumerate}[label=(\alph*)]
  \item $\Phi_8$의 $\F_3[X]$에서의 인수분해,
  \item 원시 $9$차 단위근을 포함하는 $\F_5$의 최소 확장,
  \item $X^5-1$의 $\F_7$ 위 분해체와 갈루아군.
\end{enumerate}
\end{exercise}

\subsection*{B. 핵심 정리의 증명}

\begin{exercise}
체의 곱셈군의 모든 유한부분군이 순환군임을 증명하고, 이를 이용해 $\charac K\nmid n$일 때 $\mu_n$이 위수 $n$의 순환군임을 보여라.
\end{exercise}

\begin{exercise}[*]
순환군의 원소를 정확한 위수에 따라 분류하여
\[
  \sum_{d\mid n}\varphi(d)=n
\]
을 증명하라. 이 항등식을 $n=12,18$에서 직접 검산하라.
\end{exercise}

\begin{exercise}
$p$가 소수일 때 다음 공식을 증명하라.
\begin{enumerate}[label=(\alph*)]
  \item $\Phi_{p^r}(X)=\Phi_p(X^{p^{r-1}})$,
  \item $p\nmid n$이면
  \[
    \Phi_{pn}(X)=\frac{\Phi_n(X^p)}{\Phi_n(X)},
  \]
  \item $n>1$이 홀수이면 $\Phi_{2n}(X)=\Phi_n(-X)$.
\end{enumerate}
\end{exercise}

\begin{exercise}[*]
원시 $n$차 단위근 $\zeta_n$의 최소다항식을 $f$라 하자. $p\nmid n$인 모든 소수 $p$에 대해 $\zeta_n^p$도 $f$의 근임을 법 $p$에서 $X^n-1$의 분리성을 이용해 증명하라. 이를 바탕으로 $f=\Phi_n$과 $\Phi_n\in\Z[X]$를 증명하라.
\end{exercise}

\begin{exercise}
$\gcd(m,n)=1$일 때
\[
  \Q(\zeta_m,\zeta_n)=\Q(\zeta_{mn}),
  \qquad
  \Q(\zeta_m)\cap\Q(\zeta_n)=\Q
\]
을 증명하고 갈루아군의 직접곱 분해를 구하라.
\end{exercise}

\begin{exercise}[*]
$n>2$일 때 복소켤레의 고정체가
\[
  \Q(\zeta_n)^+
  =\Q(\zeta_n+\zeta_n^{-1})
\]
임을 증명하라. 이 체가 $\Q(\zeta_n)$ 안의 모든 실수원소로 이루어지고 차수가 $\varphi(n)/2$임도 보여라.
\end{exercise}

\begin{exercise}
$L=\Q(\zeta_7)$에서
\[
  \eta=\zeta_7+\zeta_7^2+\zeta_7^4
\]
라 하자. $\eta^2+\eta+2=0$을 증명하고 $L$의 유일한 이차 중간체가 $\Q(\sqrt{-7})$임을 보여라.
\end{exercise}

\begin{exercise}[*]
$q$가 소수거듭제곱이고 $\gcd(q,n)=1$이라 하자. 프로베니우스 궤도를 사용하여 $\Phi_n$이 $\F_q[X]$에서
\[
  \frac{\varphi(n)}{\operatorname{ord}_n(q)}
\]
개의 기약인수로 분해되고, 각 인수의 차수가 $\operatorname{ord}_n(q)$임을 증명하라.
\end{exercise}

\begin{exercise}
$X^5-1\in\F_7[X]$의 분해체를 구하고, 상대 프로베니우스가 다섯 근에 유도하는 순열을 적어라.
\end{exercise}

\begin{exercise}[*]
힐베르트 정리 90을 $\C/\R$에 적용하라. $x,y\in\Q$에 대해
\[
  x^2+y^2=1
\]
일 필요충분조건이 어떤 $m,n\in\Q$에 대해
\[
  x=\frac{m^2-n^2}{m^2+n^2},
  \qquad
  y=\frac{2mn}{m^2+n^2}
\]
로 쓸 수 있는 것임을 보여라.
\end{exercise}

\subsection*{C. 순환확장과 Kummer 구조}

\begin{exercise}
$L/K$가 유한 순환 갈루아 확장이고 $\Gal(L/K)=\langle\sigma\rangle$라 하자. 체매장의 선형독립성을 이용하여
\[
  \operatorname N_{L/K}(b)=1
  \quad\Longrightarrow\quad
  b=\frac{a}{\sigma(a)}
\]
인 $a\in L^\times$가 존재함을 증명하라.
\end{exercise}

\begin{exercise}[*]
$\charac K\nmid n$이고 $\mu_n\subseteq K$라 하자. 다음을 증명하라.
\begin{enumerate}[label=(\alph*)]
  \item 모든 차수 $n$의 순환 갈루아 확장은 $X^n-c$의 한 근으로 생성된다.
  \item $a^n=c\in K^\times$이면 $K(a)/K$는 순환 갈루아 확장이고 그 차수는 $n$을 나눈다.
\end{enumerate}
\end{exercise}

\begin{exercise}[*]
$\charac K=p>0$이라 하자. 가법형 힐베르트 정리 90을 사용하여 다음 아르틴--슈라이어 정리를 증명하라.
\begin{enumerate}[label=(\alph*)]
  \item 차수 $p$의 모든 순환 갈루아 확장은 $X^p-X-c$의 한 근으로 생성된다.
  \item $X^p-X-c$는 $K$에서 완전히 분해되거나 기약이다.
  \item 기약인 경우 분해체의 갈루아군은 $C_p$이다.
\end{enumerate}
\end{exercise}

\begin{exercise}
$X^3-2$를 사용하여 다음 두 확장을 비교하라.
\[
  \Q(\sqrt[3]{2})/\Q,
  \qquad
  \Q(\zeta_3,\sqrt[3]{2})/\Q(\zeta_3).
\]
각 확장의 차수, 정규성, 분리성 및 가능한 자동동형을 구하라.
\end{exercise}

\begin{exercise}[*]
$K=\Q(\zeta_5)$, $a=\sqrt[5]{2}$라 하자. $[K(a):K]=5$임을 보이고
\[
  \Gal(K(a)/K)\cong C_5
\]
임을 증명하라. 생성원의 $a$와 $\zeta_5$에 대한 작용도 적어라.
\end{exercise}

\begin{exercise}
$L=\Q(\sqrt2,\sqrt3)$에 대해 Kummer 쌍
\[
  \Gal(L/\Q)\times
  \langle[2],[3]\rangle
  \longrightarrow\mu_2
\]
의 값을 표로 작성하라. 이 쌍이 비퇴화임을 직접 확인하라.
\end{exercise}

\begin{exercise}
다음 체에 들어 있는 모든 단위근을 구하라.
\[
  \Q(\sqrt2),\qquad
  \Q(i),\qquad
  \Q(\sqrt{-2}),\qquad
  \Q(\sqrt{-3}).
\]
힌트: $\zeta_n$이 이차체에 들어가면 $\varphi(n)\le2$이다.
\end{exercise}

\begin{exercise}
$p$가 소수이고
\[
  p-1=\ell_1\cdots\ell_r
\]
가 서로 다른 소수들의 곱이라고 하자. $\Q(\zeta_p)/\Q$의 중간체가 정확히 $2^r$개임을 보여라. 각 가능한 중간체 차수도 설명하라.
\end{exercise}

\begin{exercise}
$K=\Q(\zeta_3)$와
\[
  L=K(\sqrt[3]{2},\sqrt[3]{3})
\]
에 대해 $[L:K]=9$와
\[
  \Gal(L/K)\cong C_3\times C_3
\]
을 증명하라.
\end{exercise}

\begin{exercise}[*]
$L=\Q(\zeta_{15})$에 대해 다음을 수행하라.
\begin{enumerate}[label=(\alph*)]
  \item $\Gal(L/\Q)\cong C_4\times C_2$임을 보여라.
  \item 위수 $4$인 부분군이 몇 개인지와 이차 중간체가 몇 개인지 구하라.
  \item 세 이차 중간체가
  \[
    \Q(\sqrt{-3}),\qquad
    \Q(\sqrt5),\qquad
    \Q(\sqrt{-15})
  \]
  임을 보여라.
\end{enumerate}
\end{exercise}
